% PDF-derived LaTeX transcription. Original manuscript source was unavailable.
% Published PDF SHA-256: dae11ba9b0d07d8fc654004e1d8b048bb5cc865a6a960cdc8fd7ef80c6b1432e
% Extracted text SHA-256: b931b9492a59f3e3d5e537b08051964d527ad62ea82dfcb72627053812589ca7
\documentclass[11pt]{article}
\usepackage[margin=1in]{geometry}
\usepackage{amsmath,amssymb}
\setlength{\emergencystretch}{3em}
\title{Finite Equivalence and Nested Performance Limits in a Fused Reactor-Field Generator}
\author{}
\date{}
\begin{document}
\maketitle
\noindent\textbf{PDF-derived LaTeX transcription.} This editable text follows the PDF extraction. Refer to the published PDF for figures, table layouts, and mathematical notation.

% Extracted chunk 1

Finite Equivalence and Nested Performance Limits in a Fused
Reactor-Field Generator

\section*{Abstract}

Kernel fusion can preserve a numerical reference while changing the dominant costs of sci-
entific data generation. We examine this interaction in a bounded reduced-order stirred-reactor
field generator with FP64 transported state and FP32 stored output. The evidence distinguishes
an original Torch GPU port, a first transport/reaction fusion, and a later merged-kernel compar-
ison against a conservative shared-face thermal reference. The first fusion reduced the measured
stencil section from 48.10694 to 1.40098 ms per step, a factor of 34.34, while a matched eight-
sample campaign decreased from 428 to 111 integer shell seconds. These are different timing
scopes. Stored-FP32 parity to NumPy, finite FP64 bitwise agreement with Torch, and saved
digest agreement were also separate tests. The later merge passed 18 unit cases and three
batch-three trajectory comparisons containing 162 paired saved-field digests. On a shared node,
its device sequence improved by 1.39238 times, but the complete eight-sample chunk, including
saves, improved by only 1.10104 times. Concurrent campaigns exposed host-memory limits,
including an out-of-memory failure. An earlier 48-triple eager precision study rejected both
FP32 variants under a conjunctive energy-ledger contract despite small temperature errors and
transport defects below the specified bound; later native-FP32 timing and smoke results con-
cern a different implementation. The case demonstrates useful finite numerical equivalence and
substantial local acceleration, but also shows why neither establishes universal equality, physical
calibration, or a multiplicative production-throughput gain.

\section*{1 Introduction and contributions}

Scientific field generators often combine numerical integration with output conversion, validation,
serialization, and downstream encoding. Optimizing the integration loop can therefore change,
rather than remove, the limiting resource. A second difficulty is numerical: preserving stored data
is not the same as preserving every internal floating-point operation, and neither is evidence that
the underlying physical model is accurate.
This paper studies these distinctions in an explicit advection–diffusion–reaction and heat genera-
tor for a procedurally defined stirred vessel. Its intended role is bounded field-contract testing and
data-pipeline development. The flow is an analytic moving-reference-frame (MRF) envelope, not a
resolved rotating-mesh calculation. The reaction mechanism and thermal properties are declared
synthetic pilot inputs rather than plant-calibrated chemistry.
The focused contribution is a retrospective analysis of finite numerical equivalence and nested per-
formance limits. We distinguish three implementation stages. The original Torch port establishes
the eager GPU arithmetic baseline. The first FP64 fusion combines transport and pointwise reac-
tion/update work into separate fused kernels, together with caching, synchronization, save-path,
and host-pipeline changes. A later comparison starts from a different, conservative shared-face
1

thermal reference and merges its transport and reaction/update kernels while retaining thermal-
face work and reductions separately. These stages do not share a single unchanged numerical or
timing baseline.
We organize the evidence around three questions. First, which quantities were actually equal,
and over which finite test populations? Second, how much of a local kernel gain survives at chunk,
campaign, and concurrent-node scope? Third, which adverse outcomes constrain extensions to
greater concurrency or lower precision? This organization retains successful comparisons alongside
slow repeats, metadata differences, precision-gate failures, and host out-of-memory failures. It does
not treat implementation comments or acceptance labels as experimental proof.

\section*{2 Related work}

Triton provides tile-based programming and compiler optimization for custom GPU computation
[1]. That work supplies context for expressing fused tensor operations with a tiled compiler; the
numerical equivalence and timings examined here depend on the captured reactor implementation
and its tests, not on the compiler paper.
The Roofline technical report relates attainable floating-point performance to operational inten-
sity and sustainable memory bandwidth, and motivates bound and bottleneck analysis [ 2]. We use
this perspective only to distinguish resource limits. The cited version is University of California,
Berkeley Technical Report UCB/EECS-2008-134, not the later journal publication. No hardware-
counter traffic measurements or sustainable-bandwidth characterization are available for the present
case, so we do not construct a measured Roofline or claim bandwidth-optimal execution.
The present study is consequently narrower than a general GPU compiler evaluation or a new
reactor discretization. It examines a specific numerical chain, the finite evidence for preserving it,
and the attenuation of measured gains as the timed workload expands.

\section*{3 Scientific methods and numerical model}


\section*{3.1 Retrospective design and computational configuration}

This analysis uses captured source, reports, logs, manifests, and scalar results. No additional GPU
simulations, raw-field measurements, or token-shard replays were performed. Saved-record consis-
tency checks can establish agreement between retained results and source identities, but cannot
substitute for rerunning the original experiments. The private inventory and original evidence are
retained in Supplementary record S1.
The reported GPU environment was an NVIDIA H100 with 80 GB device memory, Torch
2.10.0+cu128, and Triton 3.6.0. Production-shaped samples used a cubic grid of 1283 points,
640 explicit steps, \ensuremath{\Delta}t = 0 .0025 s, and saving every 40 steps. Thus the horizon was 1.6 s and the
output contained 16 frames at 0.1-s intervals, without an initial-time frame. The stored layout was
time, channel, Z, Y, X, with 18 FP32 channels. The transported state comprised six concentrations,
temperature, and mixing age, held in FP64 on the reference fused paths. Analytic velocity and
pressure were computed using FP64 expressions and converted to FP32 before transport or storage.
Initial states were constructed with the NumPy reference and transferred to the GPU; the effective
random seed was the configured seed plus sample index. Samples within a chunk shared geometry
and controls but differed in initial state.
2


\section*{3.2 Geometry and analytic flow envelope}

The procedural vessel had radius R = 0 .45 m and height H = 0 .90 m, with a masked impeller
of radius Ri = 0 .15 m at height zi = 0 .45 m. The grid spacing was h = 2 R/(N \ensuremath{-} 1), where N
is the number of grid points per axis. The active-fluid mask excluded vessel and impeller solids;
the jacket mask identified sidewall exchange. Signed-distance geometry was inserted during final
packing, replacing the simulator’s placeholder geometry plane.
For cylindrical radius r, define \ensuremath{\xi} = min(r/R, 1), \ensuremath{\zeta} = z/H , and
W = (1 \ensuremath{-} \ensuremath{\xi}2)2, V = sin 2(\ensuremath{\pi}\ensuremath{\zeta} ), (1)
I = exp
[
\ensuremath{-}
( z \ensuremath{-} zi
0.22H
) 2]
. (2)
If n(t) is the impeller speed in revolutions per minute, \ensuremath{\omega}(t) = 2 \ensuremath{\pi}n(t)/60 has units s \ensuremath{-}1. The
ramp duration was tr = max(10\ensuremath{\Delta} t, 0.2Nt\ensuremath{\Delta}t), with n(t) = n\ensuremath{\infty} min(max(t/tr, 0), 1). Here Nt is the
number of explicit time steps. The prescribed envelope was
u\ensuremath{\theta} = 0.28\ensuremath{\omega}rW V (0.35 + 0.65I), (3)
a = 0.018\ensuremath{\omega}Ri, (4)
ur = \ensuremath{-}arW \ensuremath{\pi}
H sin(2\ensuremath{\pi}\ensuremath{\zeta} ), (5)
uz = 2a(1 \ensuremath{-} \ensuremath{\xi}2)(1 \ensuremath{-} 3\ensuremath{\xi}2)V. (6)
The cylindrical components were transformed to Cartesian components, with the inverse radius set
to zero on the axis. Velocity was zeroed in solids. The centrifugal pressure estimate was
p = \ensuremath{\rho}u2
\ensuremath{\theta}
2 \ensuremath{-}
\ensuremath{\langle} \ensuremath{\rho}u2
\ensuremath{\theta}
2
\ensuremath{\rangle}
fluid
, (7)
with zero solid pressure. Here the brackets denote an active-fluid spatial mean. This is a framewise
gauge estimate, not a pressure-Poisson solution. The analytic envelope and solid masking do not
establish a resolved incompressible flow around moving blades.

\section*{3.3 Transport, reactions, and stored observables}

For a transported scalar f , the cellwise operator was
LDf = D
\ensuremath{\sum}
j=x,y,z
f +
j \ensuremath{-} 2f + f \ensuremath{-}
j
h2 \ensuremath{-}
\ensuremath{\sum}
j=x,y,z
ujGjf, (8)
where Gjf = (f \ensuremath{-} f \ensuremath{-}
j )/h for nonnegative uj, and (f +
j \ensuremath{-} f )/h otherwise. Inactive or out-of-domain
neighbors were replaced by the center value. This imposed the implemented no-flux stencil without
periodic wrapping. Axis contributions were accumulated in the fixed order x, y, z. This cellwise
advective operator should not be confused with the later shared-face thermal operator.
Writing concentrations as A, B, C, P, S1, S2, the synthetic network was
A + B + C \ensuremath{-} \ensuremath{\to}P, A + B \ensuremath{-} \ensuremath{\to}S1, P \ensuremath{-} \ensuremath{\to}S2. (9)
3

The rates were rd = kdABC , r1 = k1AB, and r2 = k2P , in mol m \ensuremath{-}3 s\ensuremath{-}1. Rate constants used
kj(T ) = kj,ref exp
[
clip
(
\ensuremath{-} Ej
Rg
( 1
Ts
\ensuremath{-} 1
Tref
)
, \ensuremath{-}30, 30
)]
, (10)
where Ts = clip( T, 250, 550) K, Tref = 330 K, and Rg is the molar gas constant. In order d, 1, 2,
activation energies were 42000, 47000, 56000 J mol \ensuremath{-}1, and reference constants were 10\ensuremath{-}5 m6 mol\ensuremath{-}2
s\ensuremath{-}1, 2 \ensuremath{\times} 10\ensuremath{-}4 m3 mol\ensuremath{-}1 s\ensuremath{-}1, and 2 \ensuremath{\times} 10\ensuremath{-}3 s\ensuremath{-}1. These values define the pilot, not calibrated reaction
kinetics.
The species source vector was
s = (\ensuremath{-}rd \ensuremath{-} r1, \ensuremath{-}rd \ensuremath{-} r1, \ensuremath{-}rd, rd \ensuremath{-} r2, r1, r2). (11)
Each active concentration used the explicit update cn+1 = max(cn + \ensuremath{\Delta}t(LDcn + sn), 0), with zero
inactive state. Heat release was ˙q = \ensuremath{-} \ensuremath{\sum}
j \ensuremath{\Delta}Hjrj, with enthalpies \ensuremath{-}82000, \ensuremath{-}45000, \ensuremath{-}18000 J mol\ensuremath{-}1.
The declared properties were \ensuremath{\rho} = 1020 kg m \ensuremath{-}3, cp = 3900 J kg \ensuremath{-}1 K\ensuremath{-}1, species diffusivity 2 \ensuremath{\times} 10\ensuremath{-}5
m2 s\ensuremath{-}1, thermal diffusivity 1.4 \ensuremath{\times} 10\ensuremath{-}5 m2 s\ensuremath{-}1, and jacket coefficient 450 W m \ensuremath{-}2 K\ensuremath{-}1.
The first-fusion thermal update used cellwise transport plus ˙q/(\ensuremath{\rho}cp) and a jacket-band source
hc(Tj \ensuremath{-} T )/(\ensuremath{\rho}cph). Mixing age used the same cellwise advection with half the species diffusivity,
a unit active-fluid source, and nonnegative clipping. An upper clamp of \ensuremath{\Delta}t/2 was applied to all
active cells in the horizontal band |z \ensuremath{-} zi| \ensuremath{\leq} 1.5h, without a radial restriction to the impeller radius.
It is an impeller-contact mixing-age diagnostic, not a through-flow residence-time distribution.
Stored observables comprised velocity components, gauge pressure, temperature, the six con-
centrations, desired reaction rate, heat release, mixing age, signed distance, and three masks. The
saved state followed the update, while reaction-rate and heat-release snapshots used rates evaluated
from the old state at that step. This sampling convention was part of the chain being preserved.

\section*{3.4 Conservative shared-face thermal variant}

The later reference replaced cellwise temperature transport with shared-face exchange. For an
oriented internal fluid face from cell l to u, its implemented advective and diffusive powers were
Fa = \ensuremath{\rho}cpuf (Tup \ensuremath{-} T )h2, (12)
Fd = \ensuremath{-}\ensuremath{\rho}cpDT
Tu \ensuremath{-} Tl
h h2, (13)
where uf was the mean of adjacent stored velocity components and T was the active-fluid mean
temperature at that fine step. The upwind temperature Tup was Tl for uf \ensuremath{\geq} 0 and Tu otherwise;
both neighboring temperatures were old-state values taken before the update. Each signed face
power was subtracted from cell l and added to cell u, so positive power transferred heat from l to
u. This pairwise exchange conserves the domain sum independently of discrete velocity divergence.
It does not establish local equivalence between mean-referenced and absolute-temperature advec-
tion when the masked discrete velocity has nonzero divergence. Inactive faces carried no internal
transport; explicit fluid–jacket faces exchanged hc(T \ensuremath{-} Tj)h2 W, with outward heat positive and
face multiplicity retained. Temperature rates followed by division by \ensuremath{\rho}cph3.
Every fine step contributed to interval-integrated energies Qa, Qd, Qr, Qj, in joules. With U =
\ensuremath{\rho}cph3 \ensuremath{\sum}
fluid T , the ledger compared
\ensuremath{\Delta}Upred = Qa + Qd + Qr \ensuremath{-} Qj (14)
4

to the measured state-energy change. Its two dimensionless tests were
eb = |\ensuremath{\Delta}U \ensuremath{-} \ensuremath{\Delta}Upred|
max(|\ensuremath{\Delta}U |, |Qa| + |Qd| + |Qr| + |Qj|, 1 J) , (15)
et = |Qa| + |Qd|
max(|\ensuremath{\Delta}U |, |Qr| + |Qj|, 1 J) . (16)
Every interval required both eb \ensuremath{\leq} 10\ensuremath{-}8 and et \ensuremath{\leq} 0.02. The second quantity measures closed-
domain net transport imbalance, not spatial error against an independently known temperature
field. Legacy cellwise thermal mode carried no corresponding conservation claim.

\section*{4 Validation protocols and timing definitions}


\section*{4.1 Four distinct evidence levels}

Stored-output parity to NumPy required zero maximum absolute error over compared FP32 fields
and a passing domain audit. This is a finite comparison of stored values. FP64 bitwise gates instead
compared the fused arithmetic with the accepted eager Torch chain, including integer views of state
tensors where specified. A third level compared retained raw-field, flux, metadata, diagnostic, or
latent digests. Digest agreement concerns the serialized quantities covered by those digests; it is not
a new measurement of the original arrays. Independent physical accuracy would require external
reference solutions or experimental observations and was not tested here.
The runtime stability limits were \ensuremath{\Delta}t maxfluid
\ensuremath{\sum}
j |uj|/h \ensuremath{\leq} 0.45, \ensuremath{\Delta}t max(Ds, DT )/h2 \ensuremath{\leq} 0.16, and
maximum reaction consumption fraction per step at most 0.1. Consumption ratios used old con-
centrations with a 10\ensuremath{-}8 mol m \ensuremath{-}3 concentration floor. Reaction violations and NaNs accumulated
on the device and were checked at saves and loop completion before affected output could proceed.
Domain checks also covered finite fields, binary static masks, absolute solid-velocity components
no greater than 10\ensuremath{-}7 m s \ensuremath{-}1, absolute framewise fluid-mean gauge pressure no greater than 10\ensuremath{-}2
Pa, nonnegative concentrations and rates, fluid temperatures of 250–550 K, and the conservative
ledger when applicable. The solid-velocity audit tolerance is distinct from the model’s exact zero
assignment in solids.

\section*{4.2 Operation-order discipline and finite test populations}

The first fusion disabled fused multiply-add contraction, passed scalar parameters through FP64
tensors, and reproduced Torch CUDA scalar division as multiplication by a precomputed reciprocal.
These choices constrained rounding order rather than proving an algebraic equivalence theorem.
Its final record contains six stencil cases: grids 32 and 48, each with three seeds and both geometry-
derived and random activity patterns across the set. It also records fused-update on/off trajectory
comparisons at those two grids over 640 steps, and stored-FP32 NumPy parity at grids 32, 48
with multiple chunks, and 128. The excerpted final parity record does not enumerate all sample
counts, so these three grid configurations are not presented as a larger known trajectory population.
A separate eight-sample production-grid campaign compared raw, interface-flux, and metadata
digests.
The later merged-kernel unit gate contained 18 enumerated cases at grids 32 and 48, batch
B = 3. Cases covered conservative and legacy thermal modes, real and random states, real masks
and random masks activating boundary planes, and selected heat-output settings. This was not a
5

full Cartesian product. Full-trajectory comparisons used separate processes, seed 8800 with sample
offsets 0, 1, 2, and 640 steps: conservative mode at both grids and legacy mode at grid 32. The
resulting three batch comparisons comprised nine sample-trajectory pairs and 3\ensuremath{\times}3\ensuremath{\times}18 = 162 saved-
field digest comparisons, plus diagnostic-list digests. A further production-shaped comparison used
one 1283 sample. The harvester comparison used only the 16 overlapping sample identifiers in its
larger reference manifest.

\section*{4.3 Timing scopes}

All speed ratios below are treference/tcandidate within the stated scope. A stencil-section time is not
a solver time. The original section profiler synchronized sections, measured 20 steps, and emulated
saves every 20 steps; section sums can exceed an unsynchronized loop time. Solver/chunk wall
time includes initialization and saved-array production, but not necessarily packing, validation,
writing, or encoding. Campaign wall time includes solver and host productionization, including
interface-flux output in the matched runner comparison. Harvester time additionally includes con-
text processing, encoding, shard handling, and draining. Concurrent-node elapsed time measures
simultaneous campaigns and exposes shared host resources.
For the later device-sequence benchmark, CUDA events timed 60 steps after 8 warmup steps. The
sequence included the fluid-temperature mean, thermal-face kernel, and either the transport/update
pair or merged kernel. The old in-place and new ping-pong timing loops evolved their states
independently; their timings are not equality tests. Full-chunk timing used two calls per path, with
the second reported to reduce compilation influence. Measurements were taken on a shared node
and are indicative, not an isolated scaling study.
Campaign shell timings were integer seconds. The optimized runner also supplied an inner
floating-point campaign timer. Both conventions are retained rather than treating their small
difference as numerical measurement precision.

\section*{5 Experimental results}


\section*{5.1 Finite equivalence and pre-existing NumPy differences}

All six final stencil cases matched the eager Torch FP64 chain bitwise. The two recorded fused-
update trajectory comparisons matched stored fields and diagnostics, and all three final grid-level
NumPy parity configurations reported zero stored-FP32 error with passing audits. These results
coexist with explicitly observed FP64 NumPy/GPU differences. In the stencil tests, maximum
absolute internal differences ranged from 1.819 \ensuremath{\times} 10\ensuremath{-}12 to 2.910 \ensuremath{\times} 10\ensuremath{-}11. An operation-isolation
report attributed approximately 46\% of sampled random FP64 scalar-division results to one-unit-
in-the-last-place differences between reciprocal multiplication on CUDA and NumPy division. The
exponential probe covered 500000 points in the clamped range: Triton and Torch agreed bitwise,
whereas Torch CUDA and NumPy CPU were not bitwise identical, with maximum relative dif-
ference 2.22036 \ensuremath{\times} 10\ensuremath{-}16. Thus stored parity must not be restated as internal FP64 equality to
NumPy.
The eight-sample runner comparison reported matching raw, interface-flux, and metadata digests
for every sample. In the later conservative-reference comparison, all 18 unit cases passed. All 162
saved-field digest pairs and the three diagnostic digests agreed, including full thermal ledgers. The
6

single production-shaped sample had equal raw and metadata digests. These are finite historical
results, not evidence for arbitrary-input equality or independent physical accuracy.

\section*{5.2 First fusion: local gain and campaign gain}

The first implementation consolidated state, hoisted static masks, cached velocity/pressure at
unchanged angular speed, removed recurring host synchronizations, fused transport and reac-
tion/update work, and packed saves into a reused pinned transfer buffer. Output was written
directly into final host-array slices. Host packing, audits, Zarr writing, and interface-flux compu-
tation could overlap subsequent simulation work.
T able 1. First-fusion results at mixed measurement scopes: stencil-section time, recorded peak
device allocation, and matched eight-sample campaign shell time.
Measurement Reference First fusion
Stencil section (ms/step) 48.10694 1.40098
Recorded peak device allocation (GB) 13.88315 4.91363
Matched campaign shell wall (s) 428 111
Matched campaign samples 8 8
The stencil-section ratio in Table 1 was 34.3381. The original 32-sample solver benchmark at
chunk 8, including saves, recorded 10.17066 s/sample. An exclusive-node control recorded 9.56183
s/sample for the original and 8.78420 s/sample for the optimized path. Increasing optimized chunk
size to 16 and 32 gave 10.75913 and 16.14312 s/sample, respectively. Larger chunks did not improve
that save-including workload.
The changed-output isolation run, saving only at step 640, recorded 0.54052 s/sample. It demon-
strates the importance of output frequency, but is not a matched whole-solver speedup against the
standard save-every-40-steps baseline.
For the matched eight-sample campaign, the original integer shell time was 428 s, or 53.5
s/sample; the optimized shell time was 111 s. Their shell ratio was approximately 3.86. The op-
timized inner campaign timer was 110.713714 s, or 13.839214 s/sample. The original log supplied
solver time but no equivalent inner campaign timer. A separate optimized 32-sample campaign
recorded 293.871418 inner seconds, or 9.18348 s/sample, and 295 shell seconds. This different
batch is not a matched comparison with the eight-sample baseline. A later final-module 32-sample
campaign recorded 9.75355 inner s/sample.

\section*{5.3 Save-path variation and adverse repeats}

The variation was large enough to affect interpretation. One-chunk asynchronous saving recorded
1.42459 s/sample, while the same node’s one-chunk synchronous case recorded 7.05477. For 32
samples on that node, synchronous and asynchronous cases recorded 1.76706 and 3.93725 s/sample,
respectively. The asynchronous choice was therefore not uniformly faster across these runs. Three
final-module solver repeats elsewhere recorded 10.00102, 9.79494, and 9.27748 s/sample, each with
the same recorded peak device allocation. These slow results are part of the evidence, not outliers
removed to support the best case. Reports proposed allocator, page-touch, and transparent-huge-
page explanations, but no controlled intervention established a specific root cause.
7


\section*{5.4 Concurrent campaigns and host-memory limits}

Eight concurrent campaigns each requested 32 samples, under a 256-GB host-memory allocation
and 64 CPUs. Their controls covered equipment and operating-condition changes, training, vali-
dation, and test tasks, and a longer-horizon case of 960 steps. The three attempts in Table 2 had
different buffering and worker settings.
T able 2. Concurrent-node outcomes for eight campaigns requesting 32 samples each. Rows are
recorded manifest rows; RSS denotes peak user resident-set size, and OOM denotes an
out-of-memory failure.
Configuration Rows Wall (s) RSS (GB) Outcome
Chunk 8, overlap 192 984 257.4 OOM
Chunk 4, overlap 256 663 183.1 Metadata differs
Chunk 8, serial 256 808 197.9 Completed
The first attempt used four host-promotion and four flux workers per campaign. Two campaigns
were killed by out-of-memory events; six produced their requested rows. The chunk- 4 attempt
used two promotion and four flux workers. The final chunk- 8 attempt disabled overlapping host
promotion and used two flux workers. Reported peaks are user resident-set-size measurements, not
GPU memory consumption or independently measured cgroup peaks. The approximate two-chunk
trajectory estimate of 43 GB per campaign would require 344 GB for eight campaigns, already
exceeding the allocation before other overhead.
The retrospective report describes raw/flux equality but reaction-extent metadata differences
for chunk 4, and full identity for the serial chunk- 8 attempt. The directly retained original-solver
chunk-size control is narrower: three of eight metadata digests differed between chunk 4 and chunk
8. This establishes that chunk-sensitive metadata differences also occurred without the new fusion.
It does not independently verify every row in the wider campaign claim or uniquely establish
reduction order as the cause. No complete 256-row array/metadata comparison was independently
rehashed in this analysis.

\section*{5.5 Later merge against the conservative reference}

The merged kernel read old-state neighborhoods, retained transport terms in registers, immediately
applied the pointwise update, and wrote to a ping-pong state partner. The shared-face thermal
kernel, fluid-temperature mean, diagnostic reductions, and save path remained separate. This re-
tained their reference reduction expressions and avoided materializing the transport tensor, without
claiming that all step work became one kernel.
T able 3. Shared-node conservative-scheme measurements at 1283, B = 8: device sequence and
second save-including full chunk. Ratios are reference time divided by merged-path time;
measurements are indicative.
Scope Two kernels Merged Ratio
Device sequence (ms/step) 3.38089 2.42813 1.39238
Second full chunk (s) 13.05643 11.85828 1.10104
Second chunk (s/sample) 1.63205 1.48229 1.10104
8

The measurements in Table 3 used 1283, B = 8, and the conservative thermal scheme. First-call
chunk times were 13.22102 and 12.69773 s for the reference and merged paths. The full chunk
covered 640 steps including saves and device-to-host output, but excluded a complete campaign or
harvester. Shared-node execution limits all these timing comparisons.
A one-rank harvester trial completed two campaigns, retained and drained 16 samples with no
recorded rejections, and reported consistent shard payload and manifest checks. Raw and latent
digests agreed for all 16 overlapping identifiers; metadata and context digests also agreed. The
reference contained 256 rows, but only the overlap was compared. Its 75.6-s rank wall included
compilation and bundle loading and has no matched harvester baseline from which to infer a
throughput gain.

\section*{5.6 Separate lower-precision investigations}

The earlier eager precision study used 24 configurations with two sample offsets each, giving 48
FP64/naive-FP32/shifted-FP32 triples at production shape. For each of two control regimes, it
included eight box corners and four seeded interior draws. The hotter regime spanned initial
temperatures 328–348 K, jacket temperatures 285–305 K, and 60–200 rpm; the mixing regime used
312–328 K, 295–320 K, and the same speed range. All fused kernels were disabled for this study.
Shifted FP32 stored \ensuremath{\delta}T = T \ensuremath{-} T0, with campaign-constant T0 equal to initial temperature.
It reconstructed absolute temperature for rates and saves, shifted jacket differences consistently,
and accumulated its thermal ledger in FP64 from reconstructed quantities and FP32-computed
rates. Other transported fields remained absolute FP32. This implementation differs from both
the default FP64 fusion and the later native-FP32 candidate.
T able 4. Eager precision study: audit and ledger pass counts and maximum interval defects over
48 paired FP64/naive-FP32/shifted-FP32 triples. Each precision path has a pass-count
denominator of 48.
Metric over 48 triples FP64 Naive FP32 Shifted FP32
Audit and ledger passes 48 0 0
Maximum et 1.25 \ensuremath{\times} 10\ensuremath{-}16 4.90 \ensuremath{\times} 10\ensuremath{-}8 7.77 \ensuremath{\times} 10\ensuremath{-}9
Maximum eb 1.09 \ensuremath{\times} 10\ensuremath{-}10 0.24945 1.37 \ensuremath{\times} 10\ensuremath{-}7
As summarized in Table 4, every FP32 sample failed the 10\ensuremath{-}8 energy-roundoff requirement,
while every transport defect remained below 0.02. The aggregate audit label referred to transport,
but its implementation included the conjunctive ledger result. Calling these outcomes a two-
percent transport violation would therefore be incorrect. Metadata construction rejected both
FP32 variants in all 48 cases. Nevertheless, maximum final-frame temperature errors relative to
the FP64 stored output were only 0.07421875 K for naive FP32 and 3.05176 \ensuremath{\times} 10\ensuremath{-}5 K for shifted
FP32. Small field errors did not satisfy the separate ledger contract.
Downstream encoding remained a diagnostic experiment: rejected variants used the FP64 target-
free context. Median fractions of changed token elements were 0.49694 and 0.14846 for naive and
shifted FP32. All decoded checks were finite and within the study’s 200–700-K fluid-temperature
range. Those permissive decode checks did not override failed numerical acceptance tests.
A later native shifted-FP32 candidate had a batch-one merged-kernel median of 0.125856 ms
versus 0.233088 ms for FP64, a ratio of 1.85202. A separate three-repeat batch-eight chunk com-
9

parison recorded median sample times of 1.101469 s for FP64 and 1.021254 s for native shifted
FP32: the ratio of medians was 1.07855, with passing recorded audits. Other batch/save controls
included an asynchronous batch-eight median ratio of 0.89822, favoring FP64. These later small
passing runs do not reverse the earlier cohort’s failures, nor establish blanket FP32 certification or
a production end-to-end gain.

\section*{6 Discussion and limitations}

The main finding is the separation of numerical contracts and performance boundaries. The first
fusion preserved tested Torch arithmetic while retaining stored NumPy parity, despite documented
internal NumPy/CUDA differences. The later merge preserved tested conservative-reference state
transitions and diagnostic digests. Neither result requires claiming that NumPy and CUDA imple-
ment all operations identically.
Disabling contraction, preserving accumulation order, using FP64 scalar tensors, and repro-
ducing reciprocal multiplication explain how the implementation sought finite bitwise agreement.
They do not prove equality for arbitrary inputs, different compiler releases, all GPU hardware,
or untested mask/state combinations. Saved-field agreement also does not cover every unsaved
fine-step quantity. Diagnostic digest equality extends coverage to the serialized ledger and other
recorded diagnostics, but does not turn a numerical regression test into physical validation.
The performance results show a sequence of exposed limits. Removing redundant device passes
and launches greatly reduced a local stencil section. Saved-array production then became prominent
and variable. The combined device, save-path, and host-pipeline changes improved the matched
campaign, while concurrent runs exposed the host-memory cost of overlapping production stages.
The out-of-memory canary demonstrates that a buffering strategy suitable for one campaign can
fail when replicated across a node. Serial host promotion completed under the tested allocation,
at a longer node wall time than the chunk- 4 overlapping alternative, while retaining the original
chunk size.
Similarly, the later merged sequence gained more than its save-including chunk. The native-
FP32 kernel ratio exceeded its chunk ratio. These attenuations are consistent with work remaining
outside the optimized kernel. They do not identify a unique hardware bottleneck: there are no
counters sufficient to distinguish sustainable bandwidth, cache effects, host page management,
and other contention mechanisms quantitatively. The large slow-run variation and shared-node
measurements preclude a universal throughput estimate or tight confidence interval.
No product of the reported speedups is scientifically justified. The first stencil, eight-sample
campaign, later conservative device sequence, full chunk, native-FP32 candidate, and harvester
all differ in timed scope, numerical baseline, or implementation. The changed-save experiment
additionally changes the requested output. Even equal sample dimensions do not make those
measurements interchangeable.
Precision acceptance is also contract-specific. Shifted temperature substantially reduced stored
temperature error in the eager cohort, but did not meet its strict interval energy-roundoff bound.
That is a failure under a stated numerical criterion, not evidence of general physical invalidity.
Conversely, later native-FP32 smoke passes are finite successes under a different implementation
and acceptance context. The retained logs do not fully identify the exact source and execution
environment used for the native-FP32 benchmarks, limiting attribution of these results to a precisely
specified implementation.
10

Finally, the model itself remains bounded. Its analytic flow, masked geometry, explicit reaction
network, and declared material properties provide controlled data-generation structure. They do
not resolve moving-blade hydrodynamics, turbulence, plant chemistry, or independent calorimetry.
Original output arrays were not rehashed here, and broader campaign identity claims have weaker
retained verification coverage than the enumerated later digest pairs. Failed or interrupted host-
packing successors also do not diagnose a universal multi-rank deadlock; retained failure logs include
out-of-memory events and ambiguous version names. These limitations constrain generalization
without negating the useful finite regression results.

\section*{7 Conclusion}

This case supports a bounded conclusion: FP64 fusion preserved the compared numerical ref-
erences and substantially accelerated selected device work, while broader workloads exposed save
costs, host productionization, and memory-residency limits. The later conservative-reference merge
provides particularly explicit finite coverage through enumerated unit cases, trajectory digests, and
a production-shaped sample, but only indicative shared-node performance.
The transferable methodological lesson is to report the numerical object, finite population, output
schedule, timer, and workload boundary together. Stored parity, internal bitwise equality, serial-
ized metadata agreement, and physical accuracy answer different questions. Kernel acceleration,
full-chunk acceleration, and concurrent-node throughput likewise require separate denominators.
Keeping these distinctions intact makes successful optimization informative without obscuring its
failures or extending its claims beyond the evidence.

\section*{8 Data and Code A vailability}

The private research export contains generator source snapshots, timing and parity harnesses,
scripts for the precision study, logs, reports, scalar JSON results, and selected manifests with saved
digests. Supplementary record S1 retains the original captures, source identities, local verification
inputs and recipes, and reference records. These materials are held in the existing private research
export; no public repository, archive identifier, or reviewer-access service is established, and no
redistribution permission is implied.
The historical precision protocol names saved token arrays, and the harvester results report a
retained latent shard and payload checks. Those arrays and the original raw trajectory stores
were not replayed or rehashed in this analysis. The harvester intentionally did not persist raw
trajectories. The current export is not a complete runnable public dataset and does not provide
every array needed for independent reproduction. Access and redistribution require authorization
from the private storage holders; no grant is implied here. The two cited background publications
are publicly accessible through the publication records listed in the references.

\section*{References}

[1] P. Tillet, H. T. Kung, and D. Cox. Triton: An Intermediate Language and Compiler for Tiled
Neural Network Computations. Proceedings of the 3rd ACM SIGPLAN International Workshop
11

on Machine Learning and Programming Languages (MAPL), 10 pages (2019). https://doi.
org/10.1145/3315508.3329973.
[2] S. W. Williams, A. Waterman, and D. A. Patterson. Roofline: An Insightful Visual Performance
Model for Floating-Point Programs and Multicore Architectures. University of California, Berke-
ley Technical Report UCB/EECS-2008-134 (2008). https://www2.eecs.berkeley.edu/Pubs/
TechRpts/2008/EECS-2008-134.html.
12

\end{document}
